% Thesis Exposé — UPDATED 2026-06-24
% Compile with: pdflatex expose_final.tex

\documentclass[11pt,a4paper]{article}

\usepackage[utf8]{inputenc}
\usepackage[T1]{fontenc}
\usepackage[margin=2.5cm]{geometry}
\usepackage{setspace}
\usepackage{enumitem}
\usepackage{xcolor}
\usepackage{tabularx}
\usepackage{hyperref}
\usepackage{fancyhdr}
\usepackage{titlesec}
\usepackage{booktabs}
\usepackage{tcolorbox}
\usepackage{amssymb}
\usepackage{pgfgantt}
\usepackage{array}
\usepackage{colortbl}
\usepackage{multirow}

% Colors
\definecolor{primary}{RGB}{37,99,235}
\definecolor{heading}{RGB}{30,41,59}
\definecolor{body}{RGB}{51,65,85}
\definecolor{accent}{RGB}{220,38,38}
\definecolor{boxbg}{RGB}{248,250,252}
\definecolor{border}{RGB}{226,232,240}
\definecolor{green}{RGB}{22,163,74}
\definecolor{done}{RGB}{22,163,74}
\definecolor{inprog}{RGB}{234,179,8}
\definecolor{todo}{RGB}{37,99,235}

% Page style
\pagestyle{fancy}
\fancyhf{}
\fancyfoot[C]{\thepage}
\renewcommand{\headrulewidth}{0pt}

% Section formatting
\titleformat{\section}{\Large\bfseries\color{heading}}{}{0em}{}[\vspace{-0.3em}\rule{\textwidth}{1.5pt}\color{primary}]
\titleformat{\subsection}{\large\bfseries\color{heading}}{}{0em}{}
\titleformat{\subsubsection}{\normalsize\bfseries\color{heading}}{}{0em}{}

\begin{document}

% Title / Header Area
\begin{center}
    \vspace*{1cm}
    {\Huge\bfseries\color{heading} Master's Thesis Expos\'{e}\par}
    \vspace{0.6cm}
    {\Large\bfseries\color{primary}
        Sim-to-Real Transfer in Mobile Manipulation:\\
        Drift Accumulation Analysis for Autonomous\\
        Box Stacking on the LIMO Cobot\par}
    \vspace{1cm}
    {\large By: Omar Alobaid, MSE, 29606\par}
    \vspace{0.2cm}
    {\large Supervisor: Prof.\ Ronny Hartanto\par}
    \vspace{0.2cm}
    {\large Rhine-Waal University of Applied Sciences (HSRW)\par}
    \vspace{0.2cm}
    {\large June 2026\par}
    \vspace*{1.5cm}
\end{center}

% =====================================================================
\section{Title}

\begin{tcolorbox}[colback=boxbg, colframe=primary, arc=3pt, boxrule=1.5pt,
    left=1em, right=1em, top=0.8em, bottom=0.8em]
    \textbf{Primary Title:}\\[0.3em]
    {\large Sim-to-Real Transfer in Mobile Manipulation: Drift Accumulation
    Analysis for Autonomous Box Stacking on the LIMO Cobot}
\end{tcolorbox}

\vspace{0.5em}
\noindent\textbf{Alternative titles considered:}
\begin{itemize}[label=\textbullet, leftmargin=1.5em]
    \item Autonomous Mobile Manipulation for On-Demand Manufacturing: A Sim-to-Real Study on the LIMO Cobot
    \item Iterative Pick-and-Stack with a Compact Mobile Manipulator: Quantifying Pose Drift and Sim-to-Real Gap
    \item Coordinated Navigation and Manipulation on the Agilex LIMO Cobot: From Gazebo Simulation to Physical Deployment
\end{itemize}

\vspace{0.5em}

% =====================================================================
\section{Research Questions}

\begin{tcolorbox}[colback=boxbg, colframe=primary, arc=3pt, boxrule=1pt,
    left=1em, right=1em, top=1em, bottom=1em]
    \textbf{Primary Research Question:}\\[0.5em]
    Can a compact \textbf{differential-drive} mobile manipulator (Agilex LIMO Cobot)
    autonomously perform a multi-cycle vertical box-stacking task, and how does
    accumulated navigation and manipulation error constrain the number of achievable
    stacking cycles before grasp failure?
\end{tcolorbox}

\subsection{Sub-Questions}
\begin{enumerate}[label=\textbf{\arabic*.}, leftmargin=1.5em]
    \item \textbf{Pipeline integration.} How should the ROS\,2 navigation stack (Nav2)
    and manipulation framework (MoveIt\,2) be configured and coordinated for a
    \emph{multi-step} pick-transport-place-return cycle on a single resource-constrained
    platform?
    \item \textbf{Drift characterisation.} How does position error accumulate over
    repeated navigation cycles, and at what drift magnitude does the arm's reachable
    envelope fail to encompass the target object?
    \item \textbf{Simulation-to-reality transfer.} What is the quantitative gap between
    the simulation baseline (perfect odometry, zero drift) and real-hardware performance
    --- which subsystems transfer directly and where does the gap emerge?
\end{enumerate}

\vspace{0.5em}

% =====================================================================
\section{Objectives}
\begin{enumerate}[label=\textbf{\arabic*.}]
    \item \textbf{Pipeline Integration} --- Design and implement a complete multi-cycle
    mobile manipulation pipeline integrating Nav2 (differential-drive navigation),
    MoveIt\,2 (collision-aware arm planning), and HSV+depth-camera perception,
    coordinated by a mission state machine that executes \emph{pick $\to$ transport $\to$
    stack $\to$ return} cycles without human intervention.
    \item \textbf{Stacking Experiment} --- Implement a vertical box-stacking experiment
    in Gazebo: the robot picks boxes from a source table one by one, navigates to a
    target table, places each box on top of the previous one, and returns for the next
    box. Run $N = 20$ iterations.
    \item \textbf{Drift Quantification} --- Measure per-iteration navigation position
    error and grasp success rate. Identify the iteration $N^*$ at which drift causes
    the first grasp failure. Characterise the relationship between accumulated drift and
    task failure mode.
    \item \textbf{Sim-to-Real Evaluation} --- Deploy the validated pipeline on the
    physical LIMO Cobot. Compare simulation drift profiles (upper bound: perfect odom)
    against hardware drift profiles (real encoder odometry + AMCL) to quantify the
    sim-to-real gap and identify which assumptions the simulation makes that break on
    hardware.
\end{enumerate}

\newpage

% =====================================================================
\section{Scope and Limitations}

\subsection{In Scope}
\begin{itemize}[label=\checkmark, leftmargin=1.5em]
    \item Autonomous navigation in a controlled, static indoor environment (flat ground,
    known layout, static obstacles).
    \item \textbf{Differential-drive} kinematic model --- permits in-place rotation,
    enabling the robot to align precisely with either table before each pick/place.
    \item Multi-cycle vertical box stacking: pick from Table\,1, place on Table\,2,
    repeat for up to $N=20$ cycles.
    \item 6-DoF arm manipulation of small standardised boxes (known colour, known size).
    \item Depth-camera object detection and 3-D pose estimation (HSV + pin-hole
    de-projection).
    \item Quantitative drift analysis: per-iteration position error and grasp success.
    \item Sim-to-real comparison: Gazebo (perfect odom) vs.\ physical hardware
    (AMCL + real encoders).
\end{itemize}

\subsection{Out of Scope}
\begin{tcolorbox}[colback=boxbg, colframe=accent, arc=3pt, boxrule=1pt,
    left=1em, right=1em, top=0.5em, bottom=0.5em]
    \textbf{This thesis does NOT include:}
    \begin{itemize}[label=\textbullet, leftmargin=1em, topsep=0.2em, itemsep=0.2em]
        \item Dynamic or outdoor environments.
        \item Dynamic obstacle avoidance during manipulation.
        \item Multi-robot coordination or fleet management.
        \item Advanced or learned grasping strategies (no deep learning).
        \item Perception of unknown or arbitrary objects (predefined coloured boxes only).
        \item Ackermann steering --- this kinematic mode was prototyped and abandoned
        in favour of differential drive (which permits in-place rotation).
    \end{itemize}
\end{tcolorbox}

\subsection{Limitations (Honest Assessment)}
\begin{itemize}[label=\textbullet, leftmargin=1.5em]
    \item \textbf{Simulation odometry is perfect.} The Gazebo baseline uses
    \texttt{odometry\_source: WORLD}, which produces drift-free pose estimates. This is
    intentional --- it establishes the upper performance bound. Hardware trials will
    introduce real drift.
    \item \textbf{Depth camera blind spot.} The Orbbec DaBai has a physical minimum
    sensing distance of 0.30\,m, which is larger than the grasp distance
    ($\approx$0.21\,m). This is a known constraint treated as a sim-to-real finding
    rather than a blocker.
    \item \textbf{Arm reach limits stacking height.} The myCobot 280 has a 280\,mm
    nominal reach; for top-down grasps the effective reach is shorter
    (approx.\ 180\,mm forward). A growing stack eventually leaves the arm's reachable
    envelope. This physical limit defines the maximum achievable stack height
    independently of drift.
    \item \textbf{Weld model.} Grasping in simulation uses a rigid body weld. This
    tests the navigation, perception, and planning layers but does not validate frictional
    grasp closure on the physical gripper.
\end{itemize}

\vspace{0.5em}

% =====================================================================
\section{Overview and Context}

Autonomous mobile manipulation --- the integration of autonomous navigation with robotic
arm manipulation on a single platform --- is a core capability for on-demand
manufacturing. In smart manufacturing scenarios such as the RoboCup Smart Manufacturing
League's Embodied Autonomous Intelligence Workshop (EAI-WS)~\cite{masannek2026eaiws},
robots must navigate shared spaces, locate and grasp target objects, transport them, and
place them precisely at target locations --- a complete \emph{pick-transport-place}
cycle that must execute reliably over repeated iterations.

This thesis investigates this task on the \textbf{Agilex LIMO Cobot}: a compact
mobile manipulator combining a differential-drive mobile base (LIMO PRO, 4\,kg
payload, 1\,m/s maximum speed) with a 6-DoF robotic arm (Elephant Robotics myCobot
280, 250\,g payload, 280\,mm nominal reach), an Orbbec DaBai depth camera, an EAI
T-mini Pro LiDAR, and an NVIDIA Jetson Orin Nano compute unit. The system is built
within the \textbf{ROS\,2 Humble} ecosystem using \textbf{Nav2} for autonomous
navigation (differential-drive, AMCL localisation, SLAM Toolbox mapping),
\textbf{MoveIt\,2} for collision-aware arm motion planning, and \textbf{Gazebo
Classic} for simulation.

The experimental scenario is a \textbf{multi-cycle vertical box-stacking task}: the
robot picks boxes one at a time from a source table (Table\,1), navigates to a target
table (Table\,2), places each box on top of the previous one (building a vertical
stack), returns to Table\,1, and repeats --- up to 20 cycles. This scenario is chosen
because it \emph{exercises all subsystems under compounding error}: each cycle's
starting pose depends on the navigation accuracy of the return leg, so any odometry
drift or localisation error accumulates and eventually exceeds the arm's grasp
tolerance. The iteration at which the first grasp failure occurs is a concrete,
platform-specific metric that characterises the system's operational limits.

A central contribution is the \textbf{sim-to-real analysis}: the full 20-cycle
experiment is run in simulation (establishing an upper bound with perfect odometry),
and then on the physical hardware (exposing real drift), allowing a direct quantitative
comparison of how well simulation predicts real performance.

\newpage

% =====================================================================
\section{Completed Work (As of June 2026)}

\begin{tcolorbox}[colback=boxbg, colframe=green, arc=3pt, boxrule=1pt,
    left=1em, right=1em, top=0.8em, bottom=0.8em]
    The simulation baseline is \textbf{fully operational}. Every subsystem listed below
    has been implemented, debugged, and verified in Gazebo. 29 engineering faults were
    identified and resolved.
\end{tcolorbox}

\begin{itemize}[label=\checkmark, leftmargin=1.5em, itemsep=0.5em]
    \item \textbf{Verified platform model.} Combined Xacro model with audited physical
    parameters: wheel track 0.172\,m (confirmed from datasheet), DaBai H-FOV
    67.9$^\circ$, LiDAR $\pm120^\circ$/0.10--12\,m (forward arc; self-hit filtering),
    gripper jaw range 20--45\,mm.

    \item \textbf{Navigation pipeline.} Nav2 + SLAM Toolbox + AMCL on a differential-
    drive base. Clean occupancy map built (238$\times$358 cells, 0 noise pixels, 97.2\%
    free space). Costmap tuned for arm self-masking (obstacle min range raised to
    0.35\,m). Robot navigates autonomously to any goal in the workspace.

    \item \textbf{Perception pipeline.} HSV centroid detection + depth de-projection
    via pin-hole model. Publishes 3-D \texttt{PoseStamped} on \texttt{/box\_pose}.
    Verified: correct pose estimated within grasp range.

    \item \textbf{Arm manipulation.} MoveIt\,2 with 22-seed KDL IK and RRTConnect.
    Reachability mapped: 25.9\% overall, top-down reach contracts sharply at
    $x > 0.18$\,m. Full grasp sequence (ready $\to$ pre-grasp $\to$ descent $\to$
    close $\to$ lift) executing correctly. Calibrated: lateral offset, orientation
    tolerance 0.1\,rad, closing clearance.

    \item \textbf{Integrated nav$\to$pick$\to$backup pipeline.} Single orchestrator
    node sequences travel posture $\to$ navigate $\to$ visual dock $\to$ grasp $\to$
    backup. Travel posture implemented to suppress joint-2 pendulum oscillation during
    navigation. Live TF weld service (stale-pose bug fixed). Demonstrated in simulation.

    \item \textbf{Simulation environment hardened.} GPU PRIME offload configured for
    NVIDIA RTX 3060 (eliminates depth-camera segfault on AMD iGPU). Gazebo physics
    stabilised (removed \texttt{contact\_max\_correcting\_vel=0.1} instability).

    \item \textbf{Thesis chapters drafted.} Chapters 1, 3, 4, 5, 6 written with real
    results and detailed engineering fault catalogue (29 entries). Chapter 2 (literature
    review) is the remaining writing gap.
\end{itemize}

\vspace{0.5em}

\subsection{Remaining Engineering Work}

\begin{enumerate}[label=\textbf{\arabic*.}, leftmargin=1.5em, itemsep=0.3em]
    \item \textbf{Place stage} (currently a stub) --- implement arm motion to lower
    box onto the stack and release.
    \item \textbf{Two-table navigation} --- add Table\,2 to world, add return waypoint.
    \item \textbf{Stacking height tracking} --- increment Z-target by one box height
    per cycle.
    \item \textbf{Iteration loop + data logger} --- run 20 cycles, log per-iteration
    pick/place success, nav position error, detected box pose.
    \item \textbf{Hardware deployment} --- swap Gazebo plugins for real hardware drivers,
    run 20 cycles on physical LIMO Cobot.
\end{enumerate}

\newpage

% =====================================================================
\section{Gantt Chart: Remaining Schedule}

\textbf{Reference date: 2026-06-24. Target submission: 2026-09-15.}\\[0.5em]
All simulation subsystems are complete. Remaining work: place stage, stacking
experiment, hardware trials, analysis, and writing.

\vspace{1em}

\begin{table}[h!]
\centering
\small
\setlength{\tabcolsep}{4pt}
\begin{tabularx}{\textwidth}{>{\centering\bfseries}p{1.0cm}
    >{\raggedright}p{3.2cm}
    *{12}{>{\centering\arraybackslash}p{0.6cm}}}
\toprule
\multirow{2}{*}{\textbf{Wk}} &
\multirow{2}{*}{\textbf{Focus}} &
\multicolumn{12}{c}{\textbf{Calendar (2026)}} \\
\cmidrule(l){3-14}
& &
\scriptsize Jun &
\scriptsize Jul &
\scriptsize Jul &
\scriptsize Jul &
\scriptsize Jul &
\scriptsize Aug &
\scriptsize Aug &
\scriptsize Aug &
\scriptsize Aug &
\scriptsize Sep &
\scriptsize Sep &
\scriptsize Sep \\
& &
\scriptsize 24 &
\scriptsize 1 &
\scriptsize 8 &
\scriptsize 15 &
\scriptsize 22 &
\scriptsize 29 &
\scriptsize 5 &
\scriptsize 12 &
\scriptsize 19 &
\scriptsize 26 &
\scriptsize 2 &
\scriptsize 9 \\
\midrule
\textbf{1} & Place stage + world & \cellcolor{inprog!60}$\bullet$ & & & & & & & & & & & \\
\textbf{2} & Stacking loop + logger & & \cellcolor{inprog!60}$\bullet$ & & & & & & & & & & \\
\textbf{3} & Sim experiment (20 runs) & & & \cellcolor{inprog!60}$\bullet$ & & & & & & & & & \\
\textbf{4} & Data analysis + plots & & & & \cellcolor{inprog!60}$\bullet$ & & & & & & & & \\
\textbf{5} & Literature review (Ch.\,2) & & & & & \cellcolor{todo!50}$\bullet$ & & & & & & & \\
\textbf{6} & Hardware bringup (nav) & & & & & & \cellcolor{todo!50}$\bullet$ & & & & & & \\
\textbf{7} & Hardware stacking trials & & & & & & & \cellcolor{todo!50}$\bullet$ & & & & & \\
\textbf{8} & Sim-to-real analysis & & & & & & & & \cellcolor{todo!50}$\bullet$ & & & & \\
\textbf{9} & Ch.\,3--4 finalise & & & & & & & & & \cellcolor{todo!50}$\bullet$ & & & \\
\textbf{10} & Ch.\,5--6 results/conc. & & & & & & & & & & \cellcolor{todo!50}$\bullet$ & & \\
\textbf{11} & Full draft + supervisor & & & & & & & & & & & \cellcolor{todo!50}$\bullet$ & \\
\textbf{12} & Corrections + submit & & & & & & & & & & & & \cellcolor{todo!50}$\bullet$ \\
\bottomrule
\end{tabularx}
\caption*{
    \colorbox{done!30}{\rule{0pt}{6pt}$\bullet$}\,Done\quad
    \colorbox{inprog!60}{\rule{0pt}{6pt}$\bullet$}\,In progress (current week)\quad
    \colorbox{todo!50}{\rule{0pt}{6pt}$\bullet$}\,Planned
}
\end{table}

\vspace{0.5em}

\subsection{Milestones}
\begin{table}[h]
\centering
\small
\begin{tabularx}{\textwidth}{l c X}
\toprule
\textbf{Milestone} & \textbf{Target Date} & \textbf{Deliverable} \\
\midrule
\texttt{m1-sim-complete} & \emph{Done (Jun 2026)} & Full pick$\to$nav pipeline working in sim. \\
\texttt{m2-place-stage} & Jul 7 & Place stage + two-table nav working in sim. \\
\texttt{m3-stacking-sim} & Jul 21 & 20-iteration stacking experiment in sim; CSV data. \\
\texttt{m4-literature} & Aug 4 & Chapter\,2 (literature review) complete draft. \\
\texttt{m5-hardware} & Aug 18 & 20-iteration stacking on real LIMO Cobot; CSV data. \\
\texttt{m6-analysis} & Sep 1 & Sim-to-real drift comparison plots; Ch.\,5 results. \\
\texttt{thesis-draft-1} & Sep 8 & Complete first full thesis draft for supervisor review. \\
\texttt{thesis-submit} & Sep 15 & Final corrected thesis submitted. \\
\bottomrule
\end{tabularx}
\end{table}

\newpage

% =====================================================================
\section{Risks and Mitigation}

\begin{table}[h]
\centering
\small
\begin{tabularx}{\textwidth}{X X}
\toprule
\textbf{Risk} & \textbf{Mitigation} \\
\midrule
Stack height exceeds arm reach before 20 cycles.
& The arm reach limit is a physical finding, not a failure. Report it as the
  geometric upper bound on stack height, separate from drift-induced failure. \\
\addlinespace
Real hardware odometry drift is too large to stack even 2--3 boxes.
& This is itself a result: document the minimum drift condition, propose AMCL
  re-localisation triggers between cycles as future work. \\
\addlinespace
Depth camera blind spot ($<$0.30\,m) prevents box detection at grasp range on hardware.
& Treat as sim-to-real finding. Mitigate by approaching to a slightly larger
  standoff, using remembered pose from farther detection. \\
\addlinespace
Hardware time is limited (lab access).
& Stage deployment: navigation first, then arm. Partial hardware results are
  still valid sim-to-real data. \\
\addlinespace
Literature review delayed (currently a gap).
& Allocate Week\,5 exclusively to Ch.\,2. Use the existing annotated reading list
  in \texttt{docs/references/reading\_list.md} as starting material. \\
\bottomrule
\end{tabularx}
\end{table}

\vspace{0.5em}

% =====================================================================
\section{References}

\begin{itemize}[label=\textbullet, leftmargin=1.5em]
    \item Macenski, S.\ and Jambrecic, I., ``SLAM Toolbox: SLAM for the Dynamic
    World,'' \textit{JOSS}, 2021.
    \item Macenski, S.\ et al., ``The Marathon 2: A Navigation System,'' \textit{IROS},
    2020.
    \item Masannek, M.\ et al., ``Embodied Autonomous Intelligence for On-Demand
    Manufacturing,'' EAI-WS Proposal, 2026.
    \item Dissanayaka, S.\ et al., ``From Production Logistics to Smart Manufacturing:
    The Vision for a New RoboCup Industrial League,'' 2025.
\end{itemize}

\noindent
Full bibliography maintained in \texttt{docs/thesis/bibliography.bib}.

\end{document}
